El problema dual tiene la forma:

\begin{equation}
\begin{split}
\mbox{min. } z = 80y_1 + 60y_2 + 15y_3\\
\mbox{max. } z = - 80y_1 - 60y_2 - 15y_3\\
s.a.:\\
2y_1 + 4y_2 + y_3 \geq 3\\
4y_1 + 3y_2 + y_3 \geq 2\\
y_1,\,\,y_2,\,\,y_3\geq 0\\
\end{split}
\end{equation}

La solución del dual es una solución óptima múltiple y, aplicando el método del simplex, se puede obtener una solución óptima alternativa:

\begin{center}
\begin{tabular}{c|c|ccccc|}
 & $z$ & $y_1$ & $y_2$ & $y_3$ & $h_1$ & $h_2$\\ 
 Base 1 & 45 & -50 & 0 & 0 & -15 & 0 \\ 
\hline
$h_2$ & 1  & -2 & 1 & 0 & -1 & 1\\ 
$x_3$ & 3  & 2 & 4 & 1 & -1 & 0\\ 
\hline
Base 2  & 45 & -50 & 0 & 0 & -15 & 0 \\ 
\hline
$h_2$ & 1/4  & $\frac{-5}{2}$ & 0 & $\frac{-1}{4}$ & $\frac{-3}{4}$ & 1\\ 
$y_2$ & 3/4  & $\frac{1}{2}$ & 1 & $\frac{1}{4}$ & $\frac{-1}{4}$ & 0\\ 
\hline
\end{tabular}
\end{center}




% \begin{center}
% \begin{tabular}{c|c|ccccccc|}
%   & $z$ & $y_1$ & $y_2$ & $y_3$ & $h_1$ & $h_2$ & $a_1$ & $a_2$\\ 
% Base 1 & 45 & -50 & 0 & 0 & -15 & 0 & 15 & 0 \\ 
% \hline
% $h_2$ & 1  & -2 & 1 & 0 & -1 & 1 & 1 & -1\\ 
% $y_3$ & 3  & 2 & 4 & 1 & -1 & 0 & 1 & 0\\ 
% \hline
% Base 2 & 45 & -50 & 0 & 0 & -15 & 0 & 15 & 0 \\ 
% \hline
% $h_2$ & 1/4  & -5/2 & 0 & -1/4 & -3/4 & 1 & 3/4 & -1\\ 
% $y_2$ & 3/4  & 1/2 & 1 & 1/4 & -1/4 & 0 & 1/4 & 0\\ 
% \hline
% \end{tabular}
% \end{center}
% \begin{center}
% \begin{tabular}{c|c|ccccccc|}
 % & $z$ & $y_1$ & $y_2$ & $y_3$ & $h_1$ & $h_2$ & $a_1$ & $a_2$\\ 

% \end{tabular}
% \end{center}


Por su parte, la solución del problema primal es una solución degenerada y, aplicando lemke, se puede obtener la misma solución con una base diferente:


\begin{center}
\begin{tabular}{c|c|ccccc|}
 & $z$ & $x_1$ & $x_2$ & $h_1$ & $h_2$ & $h_3$\\ 
Base 2 & -45 & 0 & -1/4 & 0 & -3/4 & 0 \\ 
\hline
$h_1$ & 50  & 0 & 5/2 & 1 & -1/2 & 0\\ 
$x_1$ & 15  & 1 & 3/4 & 0 & 1/4 & 0\\ 
$h_3$ & 0  & 0 & 1/4 & 0 & -1/4 & 1\\ 
\hline
Base 1 & -45 & 0 & -1 & 0 & 0 & -3 \\ 
\hline
$h_1$ & 50  & 0 & 2 & 1 & 0 & -2\\ 
$x_1$ & 15  & 1 & 1 & 0 & 0 & 1\\ 
$h_2$ & 0  & 0 & -1 & 0 & 1 & -4\\ 
\hline
\end{tabular}
\end{center}


La solución con la base 1 del primal se corresponde con la base 2 del dual, y se cumple el teorema fundamental de la dualidad ($y_1^*=\pi^B_1=0$, $y_2^*=\pi^B_2=0$,  $y_3^*=\pi^B_3=3$). Igualmente, La solución con la base 2 del primal se corresponde con la base 1 del dual, y se cumple el teorema fundamental de la dualidad ($y_1^*=\pi^B_1=0$, $y_2^*=\pi^B_2=3/4$, $y_3^*=\pi^B_3=0$)


